% Integración por delta de la adenda de realizaciones físicas.
% Residencia principal: transferencia física de las tres álgebras y lector temporal.

\chapter{Regímenes cuadráticos y lectura temporal}
\label{chap:regimenes-cuadraticos-tiempo}
\index{álgebra cuadrática!tres regímenes}
\index{cristal de tiempo!realización finita HMT}
\index{tiempo!lector TRIT}

\section{Álgebras cuadráticas del TRIT y signaturas de curvatura}
\label{sec:regimenes-cuadraticos-recibidos}

La construcción y la demostración primarias de los tres regímenes residen en
el \cref{chap:trit-algebras}. Para
\(\tau\in\TRIT\), con imagen balanceada
\(\bar\tau\in\{-1,0,+1\}\), se tiene
\[
 \mathbb A_\tau=\RR[X]/(X^2+\bar\tau),
\]
y, si \(J\) es la clase de \(X\),
\[
\begin{array}{ccl}
 \tau=+1&:&J^2=-I,\quad\text{régimen elíptico},\\
 \tau=0&:&J^2=0,\quad\text{régimen parabólico},\\
 \tau=-1&:&J^2=+I,\quad\text{régimen hiperbólico}.
\end{array}
\]
Para los tres valores de \(\tau\), el cociente
\(\mathbb A_\tau=\mathbb R[X]/(X^2+\bar\tau)\) satisface exactamente las tres
identidades cuadráticas anteriores. En particular, la raíz interna de
\(-1\) se realiza como operador en la hoja elíptica; HMT no necesita negar ni
suprimir la extensión compleja. El álgebra parabólica contiene un nilpotente,
y la hiperbólica se descompone mediante los dos idempotentes
\((I\pm J)/2\).

Estas identidades clasifican álgebras y transportes. No seleccionan por sí
solas una métrica espaciotemporal ni prueban que el universo físico ocupe una
de las tres hojas.

\section{Familia métrica parametrizada del lector temporal}
\label{sec:ansatz-metrico-tres-curvaturas}

Escribamos
\[
 \mathbb A_s:=\RR[X]/(X^2+s),
 \qquad s\in\{+1,0,-1\},
\]
para la misma familia con el signo real como índice.
Para hacer precisa la interfaz posible, fijemos una escala
\(\rho_{\rm curv}>0\), de dimensión \(L^{-1}\), y
\(\theta\in\RR/2\pi\ZZ\). Para \(s\in\{+1,0,-1\}\), sea
\[
 g_{s,\rho}
 =\dd r^2+S_{s,\rho}(r)^2\,\dd\theta^2,
\]
\[
 S_{s,\rho}(r)=
 \begin{cases}
  \rho_{\rm curv}^{-1}\sin(\rho_{\rm curv}r),&s=+1,\\
  r,&s=0,\\
  \rho_{\rm curv}^{-1}\sinh(\rho_{\rm curv}r),&s=-1.
 \end{cases}
\]
El dominio radial es \(r\geq0\) para \(s=0,-1\); para \(s=+1\), la carta
polar cubre \(0\leq r\leq\pi/\rho_{\rm curv}\), con la regularidad polar
usual en los extremos. En el interior de cada carta,
\[
 K_{\rm G}
 =-\frac{S_{s,\rho}''}{S_{s,\rho}}
 =s\,\rho_{\rm curv}^{\,2}.
\]
La última igualdad es un cálculo exacto del modelo una vez fijada la
realización. La flecha
\[
 \mathbb A_s\dashrightarrow(M_s,g_{s,\rho})
\]
es, sin embargo, un \emph{ansatz} de \textsc{interfaz física}; no es una
consecuencia canónica de la relación \(J^2=-sI\).

Si \(\ell_0\in\mathcal L_L^+\) es una longitud y se declara
\[
 \rho_{\rm curv}:=\frac{\Delta_4}{\ell_0},
\]
entonces
\[
 K_{\rm G}
 =s\,\frac{\Delta_4^2}{\ell_0^2}.
\]
Esta especialización está
\textsc{demostrada con estructura de partida explícita}: la conclusión se
obtiene una vez declarados \(s,\Delta_4,\ell_0\), pero no selecciona
\(\ell_0\), no
deduce la curvatura cosmológica observada ni identifica las razones
\(\mathcal G_{\rm tor}/\mathcal I_{\rm tor}\) o
\(\mathcal G_{\rm av}/\mathcal I_{\rm av}\) con \(K_{\rm G}\).

\section{Lector temporal del TRIT sobre ({-1,0,+1}): apertura, umbral y retorno con memoria}
\label{sec:lector-temporal-trit}

Una vez elegida una parametrización ordenada de las transiciones, el lector
temporal activo es
\[
\begin{aligned}
 +1&\longmapsto
 \text{retorno, memoria retenida y pasado},\\
 0&\longmapsto
 \text{umbral de evaluación y presente},\\
 -1&\longmapsto
 \text{apertura bidireccional y futuro}.
\end{aligned}
\]
El lector está \textsc{demostrado con estructura de partida explícita}: sus
conclusiones se obtienen una vez declarados la firma TRIT, el orden de la
trayectoria y la memoria transportada. No convierte el conjunto
\(\{-1,0,+1\}\) en una ecuación de movimiento. El futuro designa el espacio
de prolongaciones compatibles; la supervivencia puede seleccionar una rama
o una multisección sin borrar la genealogía de las restantes.

La frecuencia
\[
 \omega_0=\frac{2\pi}{108\,t_0}
\]
tipa el reloj del retorno completo mediante la entropía de decisión
\(I_{\mathrm{TRIT}}=\log 3\) y el transductor de Boltzmann.
No demuestra, aislada, una fase subarmónica estable. Sí permite construir el
Hamiltoniano finito de referencia que sigue; la selección de un Hamiltoniano
físico robusto requiere condiciones adicionales.

\section{Dinámica unitaria en 108 estados: Hamiltoniano, propagador y período}
\label{sec:hamiltoniano-reloj-108}

Sea
\[
 \mathcal H_{\rm clk}:=\mathbb C^{108},
 \qquad
 \{|j\rangle:j\in\mathbb Z/108\mathbb Z\}
\]
con su producto hermítico canónico. Fijemos
\(\zeta:=\exp(2\pi\mathrm i/108)\), el desplazamiento unitario
\[
 U_{\rm clk}|j\rangle=|j+1\rangle
\]
y la base de Fourier
\[
 |\widetilde k\rangle
 :=\frac1{\sqrt{108}}\sum_{j=0}^{107}\zeta^{jk}|j\rangle,
 \qquad 0\leq k\leq107.
\]
Con \(t_0\in\mathcal L_T^+\) y
\(\hbar_{\rm int}\in\mathcal L_S^+\), definimos
\[
 \boxed{
 H_{\rm clk}
 :=\hbar_{\rm int}\omega_0
 \sum_{k=0}^{107}k\,
 |\widetilde k\rangle\langle\widetilde k|,
 \qquad
 \omega_0=\frac{2\pi}{108\,t_0}.}
 \tag{CLK1}\label{eq:hamiltoniano-reloj-108}
\]
El operador tiene tipo energético
\([H_{\rm clk}]=\mathcal L_S\mathcal L_T^{-1}=\mathcal L_E\).

\begin{theorem}[Realización unitaria exacta del calendario \(108\)]
\label{thm:reloj-unitario-108}
El operador \(H_{\rm clk}\) es autoadjunto y su propagador durante un paso
satisface
\[
 U_{\rm step}
 :=\exp\!\left(-\frac{\mathrm i}{\hbar_{\rm int}}
 H_{\rm clk}t_0\right)
 =U_{\rm clk},
 \qquad
 U_{\rm step}^{108}=I.
 \tag{CLK2}\label{eq:propagador-reloj-108}
\]
Si
\[
 Z_{\rm clk}|j\rangle=\zeta^j|j\rangle,
\]
entonces, para cualquier estado de base \(|j_0\rangle\),
\[
 \langle Z_{\rm clk}\rangle_{n}
 :=\langle j_0|U_{\rm step}^{-n}Z_{\rm clk}
 U_{\rm step}^n|j_0\rangle
 =\zeta^{j_0+n}.
 \tag{CLK3}\label{eq:respuesta-reloj-108}
\]
La sucesión posee periodo primitivo \(108\).
\end{theorem}

\begin{proof}
La expresión \eqref{eq:hamiltoniano-reloj-108} es una descomposición
espectral con valores reales, luego \(H_{\rm clk}=H_{\rm clk}^\dagger\).
Además,
\[
 U_{\rm clk}|\widetilde k\rangle
 =\zeta^{-k}|\widetilde k\rangle,
\]
mientras que la exponencial de
\(\hbar_{\rm int}\omega_0 k\) durante \(t_0\) produce
\(\exp(-2\pi\mathrm ik/108)=\zeta^{-k}\). Ambos operadores coinciden sobre
una base ortonormal. Las identidades restantes siguen de
\(U_{\rm step}^n|j_0\rangle=|j_0+n\rangle\) y de la primitividad de
\(\zeta\).
\end{proof}

La formulación variacional correspondiente, para una curva
\(\psi\in C^1([t_a,t_b],\mathcal H_{\rm clk})\), no restringida durante la
variación, es
\[
 \boxed{
 L_{\rm clk}(\psi,\dot\psi)
 =\frac{\mathrm i\hbar_{\rm int}}2
 \bigl(
 \langle\psi|\dot\psi\rangle
 -\langle\dot\psi|\psi\rangle
 \bigr)
 -\langle\psi|H_{\rm clk}|\psi\rangle.}
 \tag{CLK4}\label{eq:lagrangiano-reloj-108}
\]
El funcional de acción es
\[
 \boxed{
 S_{\rm clk}[\psi]
 :=\int_{t_a}^{t_b}
 L_{\rm clk}(\psi(t),\dot\psi(t))\,\mathrm dt,}
 \tag{CLK5}\label{eq:accion-reloj-108}
\]
con \(\delta\psi(t_a)=\delta\psi(t_b)=0\).
\index{Schrödinger!ecuación del reloj finito}
Variar \(\psi\) y \(\bar\psi\) independientemente produce la ecuación de
Schrödinger del reloj finito
\[
 \boxed{\mathrm i\hbar_{\rm int}\dot\psi=H_{\rm clk}\psi}
 \tag{CLK6}\label{eq:schrodinger-reloj-108}
\]
junto con su ecuación adjunta \cite{SchrodingerNobel}. Su propagador muestreado
en \(t=nt_0\) es \(\psi_n=U_{\rm step}^{n}\psi_0\). Como \(H_{\rm clk}\) es
autoadjunto, la norma se conserva; un dato inicial normalizado permanece
normalizado. El resultado es exacto respecto de
\(H_{\rm clk},t_0,\hbar_{\rm int}\) declarados y no selecciona un Hamiltoniano
físico ni introduce una frecuencia experimental.

El \cref{thm:reloj-unitario-108} demuestra la dinámica unitaria del calendario
respecto de la completación compleja, \(t_0\) y \(\hbar_{\rm int}\)
declarados. Las cinco condiciones dinámicas de la sección siguiente reciben
una realización constructiva en \cref{sec:hmtf2-cristal-temporal}: allí se
fijan el drive, el observable de orden, la subarmonía primitiva de periodo
\(108T_{\rm d}\) y la estabilidad espectral finita. La realización
macroscópica material constituye un reconocimiento externo distinto y no una
laguna del resultado finito.

\section{Condiciones dinámicas de realización de un cristal temporal: Hamiltoniano, respuesta subarmónica y estabilidad}
\label{sec:programa-cristal-tiempo}

Una realización impulsada debe proporcionar, como mínimo:
\begin{enumerate}
 \item un espacio de estados físico y un Hamiltoniano autoadjunto
 \(H(t)\), o un protocolo impulsado equivalente, con
 \(H(t+T_{\rm d})=H(t)\);
 \item el propagador de Floquet
 \[
  U_F=\mathcal T\exp\!\left(
  -\frac{\mathrm i}{\hbar_{\rm int}}
  \int_0^{T_{\rm d}}H(t)\,\dd t\right);
 \]
 \item un observable de orden \(\mathcal O\) y una familia de estados para
 los que exista una respuesta subarmónica primitiva de orden \(m>1\):
 \[
  \langle\mathcal O((n+m)T_{\rm d})\rangle
  =\langle\mathcal O(nT_{\rm d})\rangle
 \]
 en el régimen asintótico, sin la misma periodicidad para un divisor propio
 de \(m\);
 \item estabilidad de \(m\), de la fase y de la amplitud frente a una clase
 abierta de perturbaciones admisibles;
 \item persistencia que no sea un transitorio de tamaño finito, junto con un
 mecanismo declarado de localización, pretermalización o protección
 dinámica.
\end{enumerate}
El calendario \(27\to54\to108\) proporciona la jerarquía interna de periodos.
La realización finita de \cref{sec:hmtf2-cristal-temporal} selecciona el orden
subarmónico \(108\) en su dominio declarado, sin sustituir las cinco
construcciones anteriores ni imponer por sí sola una fase material
macroscópica.

\begin{criterion}[Cristal de tiempo HMT--MD]
\label{crit:cristal-tiempo-hmt-md}
La identificación queda refutada si no existe un Hamiltoniano o protocolo
impulsado bien tipado; si la supuesta subarmonía desaparece bajo
perturbaciones arbitrariamente pequeñas compatibles con el modelo; si el
periodo observado es sólo el periodo impuesto por el impulso; si la
oscilación es un transitorio finito; o si el lector temporal asigna historias
equivalentes a fases físicas incompatibles.
\end{criterion}

\paragraph{Dominio matemático y físico del lector temporal.}
Las tres álgebras y sus exponenciales se demuestran en el dominio algebraico;
el lector temporal se deduce de la firma, el orden y la memoria iniciales. El
cristal temporal finito HMT satisface el
\cref{crit:cristal-tiempo-hmt-md} mediante la construcción de
\cref{sec:hmtf2-cristal-temporal}. La extensión a una fase material
macroscópica exige, además, especificar el soporte físico y pertenece al
contraste experimental posterior.
